\documentclass[11pt,a4paper]{article}
\usepackage[utf8]{inputenc}
\usepackage[T1]{fontenc}
\usepackage[spanish]{babel}
\usepackage{geometry}
\usepackage{amsmath}
\usepackage{amssymb}
\usepackage{graphicx}
\usepackage{enumitem}
\usepackage{booktabs}
\usepackage{array}
\usepackage{multirow}
\usepackage{fancyhdr}
\usepackage{xcolor}
\usepackage{mdframed}
\usepackage{tikz}
\usepackage{pgfplots}
\pgfplotsset{compat=1.18}

\geometry{top=2cm, bottom=2.5cm, left=2.5cm, right=2.5cm}

\pagestyle{fancy}
\fancyhf{}
\rhead{Tecnolog\'ia Digital VI --- 1er Semestre 2026}
\lhead{Parcial --- Modelo 3}
\rfoot{P\'agina \thepage}

\definecolor{tdblue}{RGB}{0,74,135}
\definecolor{lightgray}{RGB}{245,245,245}

\newmdenv[linecolor=tdblue, linewidth=1.5pt, roundcorner=4pt, backgroundcolor=lightgray, innertopmargin=6pt, innerbottommargin=6pt]{infobox}

\newcommand{\puntaje}[1]{\textbf{[#1 pts]}}

\begin{document}

%--- Encabezado ---
\begin{center}
  {\Large\bfseries Universidad Torcuato Di Tella}\\[4pt]
  {\large Tecnolog\'ia Digital VI --- Parcial Modelo 3}\\[2pt]
  {\normalsize 1er Semestre 2026}\\[8pt]
\end{center}

\begin{infobox}
\textbf{Instrucciones:} El parcial tiene una duraci\'on de \textbf{90 minutos}. Est\'a prohibido el uso de computadoras o c\'odigo. S\'i se permite el uso de una hoja de f\'ormulas. Cada secci\'on indica el puntaje. Puntaje total: \textbf{100 puntos}.
\end{infobox}

\vspace{0.5cm}

Nombre y apellido: \underline{\hspace{8cm}} \hfill Legajo: \underline{\hspace{3cm}}

\vspace{0.8cm}

%=============================================================
\section*{Secci\'on 1: Verdadero o Falso \puntaje{24}}
%=============================================================

\noindent Indique \textbf{V} (Verdadero) o \textbf{F} (Falso) para cada afirmaci\'on. Si la afirmaci\'on es falsa, \textbf{justifique brevemente} por qu\'e (sin justificaci\'on no se otorga puntaje en los falsos).

\vspace{0.3cm}

\begin{enumerate}[label=\textbf{\arabic*.}, leftmargin=*, itemsep=18pt]

\item \underline{\hspace{1cm}} SMOTE (\textit{Synthetic Minority Over-sampling Technique}) genera nuevas muestras de la clase minoritaria copiando exactamente observaciones existentes.

\vspace{0.5cm}
\noindent\textit{Justificaci\'on (si es F):} \underline{\hspace{12cm}}

\vspace{0.2cm}
\underline{\hspace{14cm}}

\item \underline{\hspace{1cm}} En K-Means, la elecci\'on aleatoria de los centroides iniciales puede llevar a diferentes resultados en distintas ejecuciones del algoritmo con los mismos datos.

\vspace{0.5cm}
\noindent\textit{Justificaci\'on (si es F):} \underline{\hspace{12cm}}

\vspace{0.2cm}
\underline{\hspace{14cm}}

\item \underline{\hspace{1cm}} La importancia de features basada en impureza en \'arboles tiende a sobreestimar la importancia de variables con muchas categor\'ias o alta cardinalidad, por lo que una alternativa m\'as confiable es la importancia por permutaci\'on.

\vspace{0.5cm}
\noindent\textit{Justificaci\'on (si es F):} \underline{\hspace{12cm}}

\vspace{0.2cm}
\underline{\hspace{14cm}}

\item \underline{\hspace{1cm}} t-SNE preserva tanto la estructura local (vecinos cercanos) como la global (distancias entre grupos lejanos) de los datos, lo que lo hace ideal para visualizaci\'on e inferencia general.

\vspace{0.5cm}
\noindent\textit{Justificaci\'on (si es F):} \underline{\hspace{12cm}}

\vspace{0.2cm}
\underline{\hspace{14cm}}

\item \underline{\hspace{1cm}} El RMSE (Root Mean Squared Error) penaliza m\'as los errores grandes que el MAE (Mean Absolute Error), ya que eleva al cuadrado las diferencias antes de promediarlas.

\vspace{0.5cm}
\noindent\textit{Justificaci\'on (si es F):} \underline{\hspace{12cm}}

\vspace{0.2cm}
\underline{\hspace{14cm}}

\item \underline{\hspace{1cm}} Un modelo de Regresi\'on Log\'istica con combinaci\'on lineal $\beta^\top x = 0$ asigna una probabilidad de 0.5 a la muestra, coloc\'andola exactamente sobre el \textit{decision boundary}.

\vspace{0.5cm}
\noindent\textit{Justificaci\'on (si es F):} \underline{\hspace{12cm}}

\vspace{0.2cm}
\underline{\hspace{14cm}}

\item \underline{\hspace{1cm}} En clustering jer\'arquico aglomerativo con \textit{Single Linkage}, la distancia entre dos grupos se define como la distancia \textbf{m\'axima} entre cualquier par de puntos pertenecientes a grupos distintos.

\vspace{0.5cm}
\noindent\textit{Justificaci\'on (si es F):} \underline{\hspace{12cm}}

\vspace{0.2cm}
\underline{\hspace{14cm}}

\item \underline{\hspace{1cm}} Gradient Boosting puede interpretarse como un descenso de gradiente en el espacio funcional, donde cada \'arbol nuevo aproxima el gradiente negativo de la funci\'on de p\'erdida respecto a la predicci\'on actual.

\vspace{0.5cm}
\noindent\textit{Justificaci\'on (si es F):} \underline{\hspace{12cm}}

\vspace{0.2cm}
\underline{\hspace{14cm}}

\end{enumerate}

\vspace{0.5cm}

%=============================================================
\section*{Secci\'on 2: M\'ultiple Choice \puntaje{24}}
%=============================================================

\noindent Marque con un c\'irculo la \textbf{\'unica} opci\'on correcta en cada pregunta.

\vspace{0.3cm}

\begin{enumerate}[label=\textbf{\arabic*.}, leftmargin=*, itemsep=20pt]

\item ¿Cu\'al de las siguientes afirmaciones sobre el algoritmo K-Means es \textbf{incorrecta}?

\begin{enumerate}[label=\alph*)]
  \item K-Means requiere que el n\'umero de clusters $K$ sea especificado de antemano.
  \item K-Means puede encontrar clusters de forma arbitraria (c\'irculos, espirales) porque no asume ninguna forma geom\'etrica particular.
  \item K-Means es sensible a los valores at\'ipicos (\textit{outliers}) ya que los centroides son promedios.
  \item En cada iteraci\'on, K-Means reasigna cada punto al centroide m\'as cercano y luego recalcula los centroides.
\end{enumerate}

\item ¿En qu\'e situaci\'on la m\'etrica RMSLE es m\'as apropiada que el RMSE?

\begin{enumerate}[label=\alph*)]
  \item Cuando la variable objetivo es binaria (0 o 1).
  \item Cuando la variable objetivo abarca varios \'ordenes de magnitud y los errores relativos importan m\'as que los absolutos (por ejemplo, predecir ventas de productos muy baratos y muy caros en el mismo modelo).
  \item Cuando el dataset tiene clases muy desbalanceadas.
  \item Cuando se quiere penalizar m\'as los errores grandes que los peque\~nos.
\end{enumerate}

\item Un cliente de un banco fue rechazado para un cr\'edito. Cuando pide explicaci\'on, el banco responde ``el modelo lo rechaz\'o''. ¿Qu\'e principio de ML aplicado fue ignorado en este sistema?

\begin{enumerate}[label=\alph*)]
  \item El principio de regularizaci\'on, ya que el modelo deber\'ia haber sido penalizado con Lasso.
  \item El principio de interpretabilidad y explicabilidad, fundamental en contextos de impacto en personas donde se debe poder justificar las decisiones autom\'aticas.
  \item El principio de validaci\'on cruzada, ya que un simple holdout no es suficiente.
  \item El principio de escalado, ya que los datos deber\'ian haberse normalizado antes del entrenamiento.
\end{enumerate}

\item ¿Cu\'al es la principal ventaja de DBSCAN frente a K-Means?

\begin{enumerate}[label=\alph*)]
  \item DBSCAN es siempre m\'as r\'apido que K-Means en datasets grandes.
  \item DBSCAN puede encontrar clusters de forma arbitraria (no esf\'erica) y detectar autom\'aticamente puntos de ruido, sin necesidad de especificar el n\'umero de clusters de antemano.
  \item DBSCAN requiere que se especifique el n\'umero de clusters $K$ como par\'ametro de entrada.
  \item DBSCAN siempre produce mejores resultados que K-Means en problemas con clases balanceadas.
\end{enumerate}

\item ¿Qu\'e representa la \textbf{altura de las uniones} en un dendrograma de clustering jer\'arquico aglomerativo?

\begin{enumerate}[label=\alph*)]
  \item El n\'umero de muestras contenidas en cada cluster al momento de la uni\'on.
  \item La distancia (o disimilitud) entre los dos grupos que se fusionan en ese paso; cuanto m\'as alta la uni\'on, m\'as diferentes eran los grupos al unirse.
  \item El valor del \'indice de Gini del nodo padre al momento de la fusi\'on.
  \item El n\'umero de iteraciones que le tom\'o al algoritmo converger hasta ese punto.
\end{enumerate}

\item En el contexto de Gradient Boosting, ¿qu\'e representa un valor de \textit{learning rate} ($\eta$) muy peque\~no (por ejemplo, $\eta = 0.01$)?

\begin{enumerate}[label=\alph*)]
  \item Que el modelo converge r\'apidamente con pocos \'arboles, sin riesgo de overfitting.
  \item Que cada \'arbol nuevo contribuye muy poco a la predicci\'on final, requiriendo m\'as \'arboles para converger, pero generando mayor regularizaci\'on impl\'icita y mejor generalizaci\'on.
  \item Que el modelo usa solo el 1\% de las muestras en cada iteraci\'on.
  \item Que la tasa de aprendizaje y el n\'umero de \'arboles son independientes entre s\'i.
\end{enumerate}

\end{enumerate}

\vspace{0.5cm}

%=============================================================
\section*{Secci\'on 3: Interpretaci\'on de Gr\'aficos \puntaje{20}}
%=============================================================

\subsection*{Ejercicio A: Curva ROC \puntaje{10}}

Se entrenaron tres modelos para un problema de clasificaci\'on binaria (detecci\'on de enfermedad). El siguiente gr\'afico muestra sus curvas ROC.

\vspace{0.3cm}
\begin{center}
\begin{tikzpicture}
\begin{axis}[
  width=9cm, height=8cm,
  xlabel={FPR (Tasa de Falsos Positivos)},
  ylabel={TPR (Tasa de Verdaderos Positivos)},
  xmin=0, xmax=1,
  ymin=0, ymax=1,
  legend pos=south east,
  grid=major,
  grid style={line width=0.3pt, draw=gray!30},
]
% Diagonal (azar)
\addplot[black, dashed, thin] coordinates {(0,0)(1,1)};
\addlegendentry{Azar (AUC=0.5)}
% Modelo A - muy bueno
\addplot[blue, thick] coordinates {
  (0,0)(0.02,0.55)(0.05,0.78)(0.10,0.88)(0.20,0.94)(0.40,0.97)(0.70,0.99)(1,1)
};
\addlegendentry{Modelo A (AUC=0.95)}
% Modelo B - moderado
\addplot[red, thick] coordinates {
  (0,0)(0.05,0.30)(0.15,0.55)(0.30,0.70)(0.50,0.80)(0.75,0.88)(1,1)
};
\addlegendentry{Modelo B (AUC=0.74)}
% Modelo C - malo
\addplot[green!60!black, thick] coordinates {
  (0,0)(0.10,0.12)(0.30,0.30)(0.50,0.51)(0.80,0.78)(1,1)
};
\addlegendentry{Modelo C (AUC=0.52)}
\end{axis}
\end{tikzpicture}
\end{center}

\vspace{0.3cm}

\begin{enumerate}[label=\textbf{\alph*)}]
  \item ¿Cu\'al de los tres modelos elegiria y por qu\'e? ¿Qu\'e indica el AUC en t\'erminos pr\'acticos?

  \vspace{2.5cm}

  \item El Modelo C tiene AUC = 0.52. ¿Qu\'e nos dice esto sobre su rendimiento? ¿Podr\'ia ser \'util de alguna forma o es siempre descartable?

  \vspace{2.5cm}

  \item En el contexto de detecci\'on de enfermedades, si se quiere \textbf{maximizar la sensibilidad} (no dejar pasar casos positivos), ¿qu\'e zona de la curva ROC del Modelo A ser\'ia de mayor inter\'es y qu\'e consecuencia tiene esto sobre los falsos positivos?

  \vspace{2.5cm}
\end{enumerate}

\subsection*{Ejercicio B: Bias-Variance Trade-off \puntaje{10}}

El siguiente gr\'afico muestra las curvas de sesgo$^2$, varianza y error total en funci\'on de la complejidad del modelo.

\vspace{0.3cm}
\begin{center}
\begin{tikzpicture}
\begin{axis}[
  width=10cm, height=6cm,
  xlabel={Complejidad del modelo},
  ylabel={Error},
  xmin=0, xmax=10,
  ymin=0, ymax=2.2,
  legend pos=north east,
  grid=major,
  grid style={line width=0.3pt, draw=gray!30},
  xtick={0,2,4,5,6,8,10},
  xticklabels={,Baja,,\'Optima,,Alta,},
]
\addplot[blue, thick] coordinates {
  (0,2.0)(1,1.6)(2,1.2)(3,0.8)(4,0.5)(5,0.3)(6,0.2)(7,0.15)(8,0.12)(9,0.08)(10,0.05)
};
\addlegendentry{Sesgo$^2$}
\addplot[red, thick] coordinates {
  (0,0.05)(1,0.08)(2,0.12)(3,0.2)(4,0.32)(5,0.45)(6,0.62)(7,0.82)(8,1.05)(9,1.3)(10,1.6)
};
\addlegendentry{Varianza}
\addplot[green!60!black, thick, dashed] coordinates {
  (0,2.05)(1,1.68)(2,1.32)(3,1.00)(4,0.82)(5,0.75)(6,0.82)(7,0.97)(8,1.17)(9,1.38)(10,1.65)
};
\addlegendentry{Error total}
\end{axis}
\end{tikzpicture}
\end{center}

\vspace{0.3cm}

\begin{enumerate}[label=\textbf{\alph*)}]
  \item ¿En qu\'e zona del gr\'afico (baja, \'optima o alta complejidad) se encuentran los fen\'omenos de \textit{underfitting} y \textit{overfitting} respectivamente? Descr\'ibalos.

  \vspace{2.5cm}

  \item ¿Por qu\'e el error total tiene un m\'inimo en la complejidad \'optima? ¿Qu\'e sucede con el sesgo y la varianza en ese punto?

  \vspace{2.5cm}

  \item Un equipo decide siempre usar el modelo m\'as complejo disponible, argumentando que ``m\'as par\'ametros siempre dan mejores resultados''. ¿C\'omo refutar\'ia este argumento usando el gr\'afico?

  \vspace{2.5cm}
\end{enumerate}

\vspace{0.3cm}

%=============================================================
\section*{Secci\'on 4: Desarrollo \puntaje{32}}
%=============================================================

\begin{enumerate}[label=\textbf{\arabic*.}, leftmargin=*, itemsep=10pt]

\item \puntaje{10} \textbf{Clustering y Aprendizaje No Supervisado}

\begin{enumerate}[label=\textbf{\alph*)}]
  \item Explique la diferencia fundamental entre el aprendizaje supervisado y el no supervisado. ¿Por qu\'e es m\'as dif\'icil evaluar el \'exito de un m\'etodo no supervisado?
  \vspace{3.5cm}
  \item Describa el algoritmo K-Means paso a paso. ¿Cu\'al es su criterio de convergencia? ¿C\'omo afecta la elecci\'on de los centroides iniciales al resultado final y qu\'e estrategia existe para mitigarlo?
  \vspace{4cm}
  \item ¿C\'omo ayuda el ``M\'etodo del Codo'' a elegir el valor de $K$ en K-Means? ¿Cu\'al es la limitaci\'on de este m\'etodo en la pr\'actica?
  \vspace{3cm}
\end{enumerate}

\item \puntaje{10} \textbf{Comparaci\'on de M\'etricas y Selecci\'on de Modelos}

\begin{enumerate}[label=\textbf{\alph*)}]
  \item Se tiene un problema de predicci\'on de precios de propiedades (regresi\'on). El equipo debate entre usar RMSE y MAE. Explique las diferencias entre ambas m\'etricas y en qu\'e situaciones preferir\'ia una sobre la otra.
  \vspace{3.5cm}
  \item Un modelo de detecci\'on de tumores malignos tiene Precision = 0.95 y Recall = 0.40. ¿Qu\'e le preocupar\'ia de este resultado? ¿Qu\'e ajuste har\'ia al modelo y c\'omo afectar\'ia el tradeoff entre Precision y Recall?
  \vspace{3.5cm}
  \item Explique la diferencia entre los m\'etodos de selecci\'on de features \textbf{Filter}, \textbf{Wrapper} y \textbf{Embedded}, dando un ejemplo de cada uno. ¿Cu\'al es computacionalmente m\'as costoso y por qu\'e?
  \vspace{3.5cm}
\end{enumerate}

\item \puntaje{12} \textbf{Caso integral: Segmentaci\'on de clientes}

Una empresa de e-commerce quiere segmentar a sus clientes para dise\~nar campa\~nas de marketing personalizadas. Tiene un dataset con:
\begin{itemize}
  \item 200.000 clientes con las siguientes variables: monto total comprado (\'ultimos 12 meses), frecuencia de compra mensual, categor\'ias de productos preferidas (10 categor\'ias, variable categ\'orica), edad, y 50 features adicionales derivadas del comportamiento en la app.
  \item No existe una variable objetivo (etiqueta) que diga a qu\'e segmento pertenece cada cliente.
\end{itemize}

\begin{enumerate}[label=\textbf{\alph*)}]
  \item ¿En qu\'e rama del Machine Learning cae este problema? ¿Es un problema de regresi\'on o clasificaci\'on? Justifique.
  \vspace{3cm}
  \item El dataset tiene 60 features en total. Antes de aplicar un algoritmo de clustering, ¿qu\'e pasos de preprocesamiento realizar\'ia y por qu\'e? Mencione al menos \textbf{tres} pasos concretos y justifique cada uno.
  \vspace{4cm}
  \item ¿Qu\'e algoritmo(s) de clustering utilizar\'ia? Compare al menos dos opciones considerando las caracter\'isticas del dataset. ¿C\'omo validar\'ia que los clusters encontrados son \textbf{\'utiles} para el negocio, dado que no hay etiquetas verdaderas?
  \vspace{4cm}
\end{enumerate}

\end{enumerate}

\vspace{1cm}
\begin{center}
\textit{--- Fin del Parcial Modelo 3 ---}
\end{center}

\end{document}
